
\chapter{连续系统振动——梁}

\intro{本章讨论工程中最常见的弯曲振动元件——梁。与弦的张力驱动不同，梁的恢复力来自弯曲刚度 $EI$。你将学习 Euler-Bernoulli 梁的四阶运动方程、各种边界条件下的固有频率（超越方程）和振型，以及模态叠加法在梁振动中的应用。}

%%%%%%%%%%%%%%%%%%%%%%%%%%%%%%%%%%%%%%%%%%%%%%%%%%%%%%%%
\section{Euler-Bernoulli 梁的弯曲振动}

\subsection{物理模型与基本假设}

梁（beam）是同时具有质量和弯曲刚度的细长结构元件。Euler-Bernoulli 梁理论基于以下基本假设：

\begin{enumerate}
    \item \textbf{平截面假设}：变形前垂直于中线的截面，变形后仍保持平面并垂直于变形后的中线;
    \item \textbf{忽略剪切变形}：截面无横向剪切应变（$\gamma_{xy}=0$）;
    \item \textbf{忽略转动惯量}：截面绕自身中性轴的转动惯性忽略不计;
    \item \textbf{小挠度假设}：挠度 $y(x,t)$ 远小于梁长和截面高度;
    \item \textbf{线弹性材料}：应力-应变关系服从 Hooke 定律。
\end{enumerate}

\begin{remark}
Euler-Bernoulli 梁理论适用于细长梁（长细比 $L/h \gtrsim 10$）。对于深梁或高频振动，需要 Timoshenko 梁理论来考虑剪切变形和转动惯量的影响。
\end{remark}

\subsection{弯曲振动方程的推导}

考虑一等截面 Euler-Bernoulli 梁，长度 $L$，截面面积 $A$，截面惯性矩 $I$，密度 $\rho$，弹性模量 $E$。

从材料力学可知，梁的弯曲曲率与弯矩的关系为：
\[
M(x,t) = EI \frac{\partial^2 y}{\partial x^2}
\]

剪力与弯矩的关系：$V = \partial M/\partial x$。

取长度为 $dx$ 的微段，如图 \ref{fig:beam-element}。

\begin{figure}[H]\centering
\begin{tikzpicture}[vib_style/.style={thick},vib_arrow/.style={-Stealth,thick}]
\draw[vib_style] (0,0) -- (8,0) node[right] {$x$};
\draw[vib_style] (0,0) -- (0,3) node[above] {$y$};
\draw[vib_style,fill=gray!20] (3,-1) rectangle (5,1);
\draw[vib_arrow,red] (2.5,0) -- (3,0) node[midway,below] {$M$};
\draw[vib_arrow,red] (5.5,0) -- (5,0) node[midway,below] {$M+dM$};
\draw[vib_arrow,blue] (3,0.8) -- (3,1.8) node[midway,left] {$V$};
\draw[vib_arrow,blue] (5,1.8) -- (5,0.8) node[midway,right] {$V+dV$};
\draw (3,0.3) -- (3,0.2) node[below] {$x$};
\draw (5,0.3) -- (5,0.2) node[below] {$x+dx$};
\end{tikzpicture}
\caption{梁微段受力分析。$M$ 为弯矩，$V$ 为剪力}
\label{fig:beam-element}
\end{figure}

微段的竖向力平衡：
\[
V - (V + dV) - f(x,t)dx = \rho A dx \cdot \frac{\partial^2 y}{\partial t^2}
\]

其中 $f(x,t)$ 为单位长度上的横向分布力。化简：
\[
-\frac{\partial V}{\partial x} - f = \rho A \frac{\partial^2 y}{\partial t^2}
\]

力矩平衡（忽略转动惯量）：
\[
(M + dM) - M - V dx = 0 \;\Rightarrow\; V = \frac{\partial M}{\partial x}
\]

将弯矩-曲率关系和剪力-弯矩关系代入力平衡方程，得到 Euler-Bernoulli 梁的弯曲振动方程：

\begin{myformula}
\[
\frac{\partial^2}{\partial x^2}\left(EI\frac{\partial^2 y}{\partial x^2}\right) + \rho A\frac{\partial^2 y}{\partial t^2} = f(x,t)
\]
\end{myformula}

对于等截面梁（$EI$ 为常数），简化为：

\begin{myformula}
\[
EI\frac{\partial^4 y}{\partial x^4} + \rho A\frac{\partial^2 y}{\partial t^2} = 0
\]
\end{myformula}

\begin{remark}
梁的弯曲振动方程是四阶空间导数——与弦振动（二阶空间导数）有本质区别。四阶导数源于弯曲刚度 $EI$ 将曲率（二阶导数）与弯矩相关联，而弯矩的二阶导数决定了惯性力。
\end{remark}

\subsection{与弦振动的对比}

\begin{table}[H]\centering
\caption{弦振动与梁振动的对比}
\label{tab:string-vs-beam}
\begin{tabular}{lcc}
\toprule
特征 & 弦（String） & 梁（Beam） \\
\midrule
恢复力来源 & 张力 $T$ & 弯曲刚度 $EI$ \\
控制方程 & $\frac{\partial^2 y}{\partial t^2}=c^2\frac{\partial^2 y}{\partial x^2}$ & $\frac{\partial^4 y}{\partial x^4}+\frac{\rho A}{EI}\frac{\partial^2 y}{\partial t^2}=0$ \\
空间导数阶数 & 二阶 & 四阶 \\
每端边界条件数 & 1个 & 2个 \\
固有频率 & $\omega_n \propto n$ & $\omega_n \propto n^2$（简支梁） \\
波速与频率关系 & 无频散（$c$ 为常数） & 有频散（$c$ 随频率变化） \\
\bottomrule
\end{tabular}
\end{table}

梁的固有频率与 $n^2$ 成正比，这意味着高阶模态的频率增长比弦更快，泛音不再是基频的整数倍（非谐波关系）。

--------------------------------------------------------------------
\section{分离变量法应用于梁}

\subsection{空间方程与时间方程}

与弦类似，假设解为 $y(x,t) = Y(x) \cdot T(t)$，代入自由振动方程：

\[
EI\,Y''''(x)T(t) + \rho A\,Y(x)\ddot{T}(t) = 0
\]

分离变量：
\[
\frac{Y''''(x)}{Y(x)} = -\frac{\rho A}{EI}\frac{\ddot{T}(t)}{T(t)} = \beta^4 \quad (\text{分离常数})
\]

注意这里取 $\beta^4$ 而非 $\beta^2$ 的原因是：时间方程的分离常数 $\omega^2$ 满足 $\omega^2 = \beta^4 \cdot EI/(\rho A)$，且 $Y''''$ 项为正。

\begin{myformula}
\[\begin{aligned}
\text{空间方程：}\quad & Y''''(x) - \beta^4 Y(x) = 0, \quad \beta^4 = \frac{\rho A \omega^2}{EI} \\
\text{时间方程：}\quad & \ddot{T}(t) + \omega^2 T(t) = 0
\end{aligned}\]
\end{myformula}

时间方程的解仍为简谐振动：$T(t) = C\sin\omega t + D\cos\omega t$。

\subsection{空间方程的通解}

空间方程 $Y'''' - \beta^4 Y = 0$ 是常系数线性 ODE。设解为 $Y = e^{\lambda x}$，特征方程为 $\lambda^4 - \beta^4 = 0$，四个根为 $\lambda = \pm\beta$（实根）和 $\lambda = \pm i\beta$（虚根）。

因此通解可写为四种等价形式之一，最常用的是：

\begin{myformula}
\[
Y(x) = C_1\sin(\beta x) + C_2\cos(\beta x) + C_3\sinh(\beta x) + C_4\cosh(\beta x)
\]
\end{myformula}

其中 $\sinh$ 和 $\cosh$ 为双曲正弦和双曲余弦。前两项对应简谐分量（与弦的解类似），后两项对应双曲分量（梁特有的，源于弯曲刚度）。

\subsection{固有频率的表达式}

由 $\beta^4 = \rho A\omega^2/EI$ 得固有频率与 $\beta$ 的关系：

\begin{myformula}
\[
\omega = \beta^2\sqrt{\frac{EI}{\rho A}}, \quad f = \frac{\beta^2}{2\pi}\sqrt{\frac{EI}{\rho A}}
\]
\end{myformula}

$\beta L$ 的值由边界条件决定的频率方程（通常是超越方程）确定。引入无量纲频率参数 $\lambda_n = \beta_n L$：

\[
\omega_n = \frac{\lambda_n^2}{L^2}\sqrt{\frac{EI}{\rho A}}
\]

--------------------------------------------------------------------
\section{典型边界条件与固有频率}

梁的每端有两个边界条件，需用几何量（挠度 $Y$、转角 $Y'$）和自然量（弯矩 $M = EI Y''$、剪力 $V = EI Y'''$）来表述。

\begin{mydefinition}{梁的边界条件类型}
\begin{itemize}
    \item \textbf{简支（铰支）}：$Y = 0$（挠度为零），$Y'' = 0$（弯矩为零）;
    \item \textbf{固支}：$Y = 0$（挠度为零），$Y' = 0$（转角为零）;
    \item \textbf{自由端}：$Y'' = 0$（弯矩为零），$Y''' = 0$（剪力为零）;
    \item \textbf{滑动端}：$Y' = 0$（转角为零），$Y''' = 0$（剪力为零）。
\end{itemize}
\end{mydefinition}

\subsection{简支梁（两端铰支）}

边界条件：$Y(0)=Y''(0)=0$，$Y(L)=Y''(L)=0$。

由 $Y(0)=0$，$Y''(0)=0$ 代入通解求得 $C_2=C_4=0$。因此 $Y(x)=C_1\sin(\beta x)+C_3\sinh(\beta x)$。

由 $Y(L)=0$，$Y''(L)=0$：
\[
C_1\sin(\beta L) + C_3\sinh(\beta L) = 0, \quad -C_1\sin(\beta L) + C_3\sinh(\beta L) = 0
\]

两式相加得 $C_3\sinh(\beta L)=0$，因 $\sinh(\beta L)\neq 0$（$\beta L>0$），故 $C_3=0$。再由 $C_1\sin(\beta L)=0$，$C_1\neq 0$，得频率方程：

\[
\sin(\beta L) = 0 \;\Rightarrow\; \beta_n L = n\pi, \quad n=1,2,3,\dots
\]

简支梁的固有频率和振型：

\begin{myformula}
\[
\omega_n = \left(\frac{n\pi}{L}\right)^2\sqrt{\frac{EI}{\rho A}} = \frac{(n\pi)^2}{L^2}\sqrt{\frac{EI}{\rho A}}, \quad Y_n(x) = \sin\left(\frac{n\pi x}{L}\right)
\]
\end{myformula}

\begin{remark}
简支梁是唯一具有与弦相同振型形式（正弦函数）的梁。这是因为简支边界条件消除了双曲分量。但频率与 $n^2$ 成正比，而非弦的 $n$。
\end{remark}

\begin{figure}[H]\centering
\begin{tikzpicture}[scale=1.3,vib_style/.style={thick},vib_arrow/.style={-Stealth,thick}]
\draw[vib_style] (0,-1.5) -- (0,1.5) node[above] {$y$};
\draw[vib_style] (0,0) -- (5,0) node[right] {$x$};
\fill (0,0) circle (0.08);
\fill (5,0) circle (0.08);
\node at (0,-0.2) {$0$};
\node at (5,-0.2) {$L$};

\draw[domain=0:5,smooth,variable=\x,blue,ultra thick] plot ({\x},{1.3*sin(pi*\x/5 r)});
\node[blue] at (5.5,1.3) {$n=1$};
\draw[domain=0:5,smooth,variable=\x,red,ultra thick] plot ({\x},{1.1*sin(2*pi*\x/5 r)});
\node[red] at (5.5,1.1) {$n=2$};
\draw[domain=0:5,smooth,variable=\x,green!60!black,ultra thick] plot ({\x},{0.9*sin(3*pi*\x/5 r)});
\node[green!60!black] at (5.5,0.9) {$n=3$};

\node at (2.5,-1.0) {简支-简支梁的前三阶弯曲振型};
\end{tikzpicture}
\caption{简支梁的前三阶弯曲振型。与弦的振型形式相同（正弦函数），但频率与 $n^2$ 成正比}
\label{fig:beam-ss-modes}
\end{figure}

\subsection{悬臂梁（一端固支一端自由）}

边界条件：$Y(0)=Y'(0)=0$（固支端），$Y''(L)=Y'''(L)=0$（自由端）。

由固支端条件可表示为 $C_2=-C_4$, $C_1=-C_3$。代入自由端条件 $Y''(L)=0$ 和 $Y'''(L)=0$ 后得到关于 $C_1, C_3$ 的齐次线性方程组，其有非零解的条件是系数行列式为零：

\begin{myformula}
\[
\cos(\beta L)\cosh(\beta L) + 1 = 0
\]
\end{myformula}

这是一个\textbf{超越方程}（transcendental equation），无法显式求解，需通过数值求根。前四阶 $\lambda_n = \beta_n L$ 的解为：

\[
\lambda_1 = 1.8751, \quad \lambda_2 = 4.6941, \quad \lambda_3 = 7.8548, \quad \lambda_4 = 10.996
\]

对于 $n > 4$，有近似公式 $\lambda_n \approx (2n-1)\pi/2$。

悬臂梁的固有频率：

\[
\omega_1 = 3.516\sqrt{\frac{EI}{\rho A L^4}}, \quad \omega_2 = 22.03\sqrt{\frac{EI}{\rho A L^4}}, \quad \omega_3 = 61.70\sqrt{\frac{EI}{\rho A L^4}}
\]

固有振型：
\[
Y_n(x) = (\cosh\beta_n x - \cos\beta_n x) - \frac{\sinh\beta_n L - \sin\beta_n L}{\cosh\beta_n L + \cos\beta_n L}(\sinh\beta_n x - \sin\beta_n x)
\]

\begin{figure}[H]\centering
\begin{tikzpicture}[scale=1.3,vib_style/.style={thick},vib_arrow/.style={-Stealth,thick}]
\draw[vib_style] (0,-1.5) -- (0,1.5) node[above] {$y$};
\draw[vib_style] (0,0) -- (5,0) node[right] {$x$};
\fill (0,0) circle (0.08);
\node at (0,-0.2) {$0$};
\node at (5,-0.2) {$L$};

\draw[domain=0:5,smooth,variable=\x,blue,ultra thick] plot ({\x},{1.3*((cosh(1.8751*\x/5)-cos(1.8751*\x/5 r)) - 0.7341*(sinh(1.8751*\x/5)-sin(1.8751*\x/5 r)))});
\node[blue] at (5.5,1.3) {$n=1$};
\draw[domain=0:5,smooth,variable=\x,red,ultra thick] plot ({\x},{0.55*((cosh(4.6941*\x/5)-cos(4.6941*\x/5 r)) - 1.0185*(sinh(4.6941*\x/5)-sin(4.6941*\x/5 r)))});
\node[red] at (5.5,0.55) {$n=2$};
\draw[domain=0:5,smooth,variable=\x,green!60!black,ultra thick] plot ({\x},{0.35*((cosh(7.8548*\x/5)-cos(7.8548*\x/5 r)) - 0.9992*(sinh(7.8548*\x/5)-sin(7.8548*\x/5 r)))});
\node[green!60!black] at (5.5,0.35) {$n=3$};

\node at (2.5,-1.0) {悬臂梁的前三阶弯曲振型};
\end{tikzpicture}
\caption{悬臂梁的前三阶弯曲振型。固支端挠度和转角为零，自由端弯矩和剪力为零}
\label{fig:beam-cantilever-modes}
\end{figure}

\subsection{两端固支梁}

边界条件：$Y(0)=Y'(0)=0$, $Y(L)=Y'(L)=0$。

频率方程：
\[
\cos(\beta L)\cosh(\beta L) - 1 = 0
\]

前四阶解：$\lambda_1=4.7300$, $\lambda_2=7.8532$, $\lambda_3=10.996$, $\lambda_4=14.137$。

近似公式（$n>2$）：$\lambda_n \approx (2n+1)\pi/2$。

\subsection{固支-简支梁}

边界条件：固支端 $Y(0)=Y'(0)=0$，简支端 $Y(L)=Y''(L)=0$。

频率方程：
\[
\tan(\beta L) - \tanh(\beta L) = 0
\]

前四阶解：$\lambda_1=3.9266$, $\lambda_2=7.0686$, $\lambda_3=10.210$, $\lambda_4=13.352$。

\subsection{典型梁的固有频率汇总}

\begin{table}[H]\centering
\caption{典型等截面梁的固有频率参数 $\lambda_n = \beta_n L$}
\label{tab:beam-frequencies}
\begin{tabular}{lcccc}
\toprule
边界条件 & $\lambda_1$ & $\lambda_2$ & $\lambda_3$ & $\lambda_n\,(n>3)$ \\
\midrule
简支-简支 & $\pi$ & $2\pi$ & $3\pi$ & $n\pi$ \\
固支-自由 & 1.8751 & 4.6941 & 7.8548 & $\approx (2n-1)\pi/2$ \\
固支-固支 & 4.7300 & 7.8532 & 10.996 & $\approx (2n+1)\pi/2$ \\
固支-简支 & 3.9266 & 7.0686 & 10.210 & $\approx (4n+1)\pi/4$ \\
自由-自由 & 0（刚体） & 4.7300 & 7.8532 & $\approx (2n-3)\pi/2$ \\
\bottomrule
\end{tabular}
\end{table}

固有频率通式：
\[
\omega_n = \frac{\lambda_n^2}{L^2}\sqrt{\frac{EI}{\rho A}}, \quad f_n = \frac{\lambda_n^2}{2\pi L^2}\sqrt{\frac{EI}{\rho A}}
\]

\begin{remark}
从表中可以看出：(1) 约束越强，固有频率越高——两端固支梁的基频是简支梁的约 2.27 倍（$4.73^2/\pi^2 = 2.27$）;(2) 自由-自由梁存在零频率的刚体模态;(3) $\lambda_n$ 随模态阶数近似等差增长，从而频率随 $n^2$ 增长。
\end{remark}

\begin{figure}[H]\centering
\begin{tikzpicture}[scale=1.2,vib_style/.style={thick},vib_arrow/.style={-Stealth,thick}]
\draw[vib_style] (0,-1.2) -- (0,1.5) node[above] {$y$};
\draw[vib_style] (0,0) -- (5,0) node[right] {$x$};
\fill (0,0) circle (0.08);
\fill (5,0) circle (0.08);
\node at (0,-0.2) {$0$};
\node at (5,-0.2) {$L$};

\draw[domain=0:5,smooth,variable=\x,blue,ultra thick] plot ({\x},{1.3*((cosh(4.7300*\x/5)-cos(4.7300*\x/5 r)) - 0.9825*(sinh(4.7300*\x/5)-sin(4.7300*\x/5 r)))});
\node[blue] at (5.5,1.3) {$n=1$};

\draw[domain=0:5,smooth,variable=\x,red,ultra thick] plot ({\x},{0.55*((cosh(7.8532*\x/5)-cos(7.8532*\x/5 r)) - 1.0008*(sinh(7.8532*\x/5)-sin(7.8532*\x/5 r)))});
\node[red] at (5.5,0.55) {$n=2$};

\node at (2.5,-0.8) {两端固支梁的前两阶弯曲振型};
\end{tikzpicture}
\caption{两端固支梁的前两阶弯曲振型。两端挠度和转角均为零}
\label{fig:beam-cc-modes}
\end{figure}

--------------------------------------------------------------------
\section{振型的正交性}

梁的振型函数满足两种正交关系：

\begin{theorem}[梁振型的正交性]
对于具有齐次边界条件的等截面梁，不同阶振型满足：
\begin{align*}
\int_0^L Y_m(x) Y_n(x)\,dx &= 0, \quad m\neq n \\
\int_0^L Y_m''(x) Y_n''(x)\,dx &= 0, \quad m\neq n
\end{align*}
\end{theorem}

\begin{proof}
$Y_n$ 满足 $EI Y_n'''' = \omega_n^2 \rho A Y_n$。两端乘 $Y_m$ 并在 $[0,L]$ 上分部积分：
\[
EI\int_0^L Y_m Y_n''''\,dx = EI\left[Y_m Y_n'''\right]_0^L - EI\left[Y_m' Y_n''\right]_0^L + EI\int_0^L Y_m'' Y_n''\,dx
\]

对于齐次边界条件，边界项为零。故：
\[
\int_0^L Y_m'' Y_n''\,dx = \frac{\omega_n^2 \rho A}{EI}\int_0^L Y_m Y_n\,dx
\]

交换 $m$ 和 $n$：
\[
\int_0^L Y_n'' Y_m''\,dx = \frac{\omega_m^2 \rho A}{EI}\int_0^L Y_n Y_m\,dx
\]

两式相减：
\[
(\omega_n^2 - \omega_m^2)\int_0^L Y_m Y_n\,dx = 0
\]

当 $m\neq n$ 时 $\omega_m\neq\omega_n$，故 $\int_0^L Y_m Y_n\,dx = 0$，进而 $\int_0^L Y_m'' Y_n''\,dx = 0$。
\end{proof}

\begin{remark}
第一正交性（质量正交性）的含义是：不同模态的惯性力互不做功。第二正交性（刚度正交性）的含义是：不同模态的弹性力互不做功。这使得模态叠加法成为可能。
\end{remark}

--------------------------------------------------------------------
\section{梁的强迫振动——模态叠加法}

\subsection{模态方程}

将受迫振动方程中的解展开为振型级数：
\[
y(x,t) = \sum_{n=1}^{\infty} Y_n(x)\,q_n(t)
\]

代入 $EI\frac{\partial^4 y}{\partial x^4} + \rho A\frac{\partial^2 y}{\partial t^2} = f(x,t)$，乘以 $Y_m$ 并积分，利用正交性：

\begin{myformula}
\[
M_n \ddot{q}_n(t) + K_n q_n(t) = Q_n(t)
\]
\end{myformula}

其中：
\[
M_n = \int_0^L \rho A\,Y_n^2\,dx, \quad K_n = \int_0^L EI\,(Y_n'')^2\,dx = \omega_n^2 M_n, \quad Q_n(t) = \int_0^L f(x,t)Y_n\,dx
\]

含阻尼的形式：
\[
\ddot{q}_n + 2\zeta_n\omega_n \dot{q}_n + \omega_n^2 q_n = \frac{Q_n(t)}{M_n}
\]

\subsection{移动载荷问题}

移动载荷是工程中重要的受迫振动问题——如车辆通过桥梁。设大小为 $F_0$ 的集中力以速度 $v$ 在梁上移动：

\[
f(x,t) = F_0\,\delta(x - vt), \quad 0 \leq vt \leq L
\]

模态力：
\[
Q_n(t) = F_0 Y_n(vt)
\]

模态方程的稳态解可计算梁在移动载荷下的动力放大因子。

\begin{remark}
移动载荷问题的特征在于：即使载荷大小恒定，由于载荷位置随时间变化，系统受到的是参数激励。当移动速度足够大时，高阶模态可能被显著激发。
\end{remark}

--------------------------------------------------------------------
\section{Timoshenko 梁简介}

Euler-Bernoulli 梁忽略剪切变形和转动惯量，当梁的截面高度与长度之比不满足细长条件（$L/h < 10$）时，或当振动频率很高时，这两项效应变得显著。

\textbf{Timoshenko 梁理论}同时考虑：
\begin{itemize}
    \item \textbf{剪切变形}：截面不再垂直于变形后的中线，存在剪切角 $\gamma$;
    \item \textbf{转动惯量}：截面绕中性轴的转动惯性 $\rho I$ 不可忽略。
\end{itemize}

Timoshenko 梁的运动方程（自由振动）：

\begin{myformula}
\[\begin{aligned}
& EI\frac{\partial^4 y}{\partial x^4} + \rho A\frac{\partial^2 y}{\partial t^2} - \rho I\left(1+\frac{E}{kG}\right)\frac{\partial^4 y}{\partial x^2\partial t^2} + \frac{\rho^2 I}{kG}\frac{\partial^4 y}{\partial t^4} = 0
\end{aligned}\]
\end{myformula}

其中 $k$ 为剪切修正因子（矩形截面 $k\approx 5/6$，圆形截面 $k\approx 9/10$），$G=E/[2(1+\nu)]$。

与 Euler-Bernoulli 梁的关键区别：
\begin{itemize}
    \item Timoshenko 梁的固有频率 \textbf{低于} 对应 Euler-Bernoulli 梁的固有频率（剪切变形和转动惯量使梁变得更``柔''）;
    \item 高阶模态的频率降幅更为显著（因剪切效应随模态阶数增大）;
    \item Timoshenko 梁存在两个频谱——弯曲波谱和剪切波谱;
    \item 对于细长梁的低阶模态，两种理论的差异很小（$<1\%$），工程中常以 $L/h > 10$ 作为 Euler-Bernoulli 理论的适用判据。
\end{itemize}

--------------------------------------------------------------------
\section{数值模拟与代码实现}

\subsection{超越方程的数值求根}

\begin{mycode}{梁的固有频率求解——超越方程的 Newton 法}
\begin{lstlisting}[language=Python]
import numpy as np
from scipy.optimize import fsolve
import matplotlib.pyplot as plt

def freq_eq_clamped_free(betaL):
    return np.cos(betaL) * np.cosh(betaL) + 1

def freq_eq_clamped_clamped(betaL):
    return np.cos(betaL) * np.cosh(betaL) - 1

def freq_eq_pinned_pinned(betaL):
    return np.sin(betaL)

def freq_eq_clamped_pinned(betaL):
    return np.tan(betaL) - np.tanh(betaL)

def find_eigenvalues(freq_eq, n_modes=5, tol=1e-10):
    lambdas = []
    search_start = 0.1
    for _ in range(n_modes):
        root = fsolve(freq_eq, search_start, xtol=tol)[0]
        lambdas.append(root)
        search_start = root + 1.5 * np.pi
    return np.array(lambdas)

print("简支梁 lambda_n:", np.round(find_eigenvalues(freq_eq_pinned_pinned), 4))
print("悬臂梁 lambda_n:", np.round(find_eigenvalues(freq_eq_clamped_free), 4))
print("固支梁 lambda_n:", np.round(find_eigenvalues(freq_eq_clamped_clamped), 4))
print("固支-简支 lambda_n:", np.round(find_eigenvalues(freq_eq_clamped_pinned), 4))
\end{lstlisting}
\end{mycode}

\begin{mycode}{绘制梁的固有振型}
\begin{lstlisting}[language=Python]
import numpy as np
import matplotlib.pyplot as plt

def cantilever_mode(beta, x):
    A = (np.cos(beta) + np.cosh(beta)) / (np.sin(beta) + np.sinh(beta))
    return (np.cosh(beta * x) - np.cos(beta * x)
            - A * (np.sinh(beta * x) - np.sin(beta * x)))

L = 1.0
x = np.linspace(0, L, 200)
lambdas = np.array([1.8751, 4.6941, 7.8548, 10.996, 14.137])

fig, ax = plt.subplots(figsize=(8, 6))
colors = ['blue', 'red', 'green', 'orange', 'purple']

for i, lam in enumerate(lambdas):
    mode = cantilever_mode(lam, x)
    mode = mode / np.max(np.abs(mode))
    ax.plot(x / L, mode + 0.3 * i, color=colors[i], lw=2,
            label=f'n={i+1}, $\\lambda_{{{i+1}}}$={lam:.4f}')

ax.set_xlabel('x / L'); ax.set_ylabel('Normalized mode shape')
ax.set_title('悬臂梁的前五阶弯曲振型')
ax.legend(); ax.grid(True, alpha=0.3)
plt.tight_layout(); plt.show()
\end{lstlisting}
\end{mycode}

\begin{mycode}{梁振动的时域响应——模态叠加法}
\begin{lstlisting}[language=Python]
import numpy as np
import matplotlib.pyplot as plt

L = 1.0
EI, rhoA = 1e4, 1.0
n_modes = 10
zeta = 0.01
t = np.linspace(0, 2, 1000)
x_obs = L / 2

lambdas = np.array([np.pi, 2*np.pi, 3*np.pi, 4*np.pi, 5*np.pi])
omega_n = lambdas**2 * np.sqrt(EI / (rhoA * L**4))

F0 = 100.0
Omega = omega_n[0] * 0.8

def mode_shape_simple(n, x, L):
    return np.sin(n * np.pi * x / L)

M_n = np.array([rhoA * L / 2] * n_modes)
f_n = F0 * np.array([mode_shape_simple(n+1, L/3, L) for n in range(n_modes)])

y_steady = np.zeros_like(t)
for n in range(n_modes):
    Y_n = mode_shape_simple(n+1, x_obs, L)
    beta = Omega / omega_n[n]
    denom = np.sqrt((1 - beta**2)**2 + (2 * zeta * beta)**2)
    phase = np.arctan2(2 * zeta * beta, 1 - beta**2)
    y_steady += Y_n * f_n[n] / (K_n := omega_n[n]**2 * M_n[n]) / denom * np.sin(Omega * t - phase)

plt.figure(figsize=(10, 4))
plt.plot(t, y_steady, 'b-', lw=1.5)
plt.xlabel('Time (s)')
plt.ylabel('Deflection at mid-span')
plt.title('简支梁在简谐集中力下的稳态响应')
plt.grid(True, alpha=0.3)
plt.tight_layout(); plt.show()
\end{lstlisting}
\end{mycode}

--------------------------------------------------------------------
\section{小结}

本章系统讨论了 Euler-Bernoulli 梁的弯曲振动：

\begin{enumerate}
    \item 梁的恢复力来自弯曲刚度 $EI$，运动方程为四阶空间导数 PDE，与弦的二阶方程有本质区别;
    \item 每端有两个边界条件，不同边界组合给出不同的固有频率和振型——频率方程多为超越方程;
    \item 简支梁是唯一具有正弦振型的梁结构;悬臂梁、固支梁的振型包含双曲函数分量;
    \item 固有频率与 $\lambda_n^2/L^2$ 成正比，高阶频率以 $n^2$ 的速度增长（非谐波关系）;
    \item 模态叠加法同样适用于梁的强迫振动分析，移动载荷是重要的工程问题;
    \item Timoshenko 梁理论修正了剪切变形和转动惯量的影响，在高频和深梁情况下不可忽略。
\end{enumerate}

